\chapter{Cystatin C}

\section*{What Is This Test?}
Cystatin C is a small protein produced by most nucleated cells and filtered by the kidney. Blood cystatin C can be used to estimate glomerular filtration rate, either alone or combined with creatinine, and may improve assessment when creatinine is affected by muscle mass.

\section*{Alternative Names}
Serum Cystatin C; Cystatin-C-Based eGFR; Kidney Filtration Marker.

\section*{Why This Test Is Ordered}
It is useful when creatinine-based eGFR may be misleading, including extremes of muscle mass, amputation, severe malnutrition, some neuromuscular disorders, or uncertainty near a kidney-function decision threshold. It may also help confirm chronic kidney disease when a creatinine-based estimate is borderline.

\section*{How This Test Is Performed}
A blood sample is collected and cystatin C is measured by an immunoassay. A laboratory calculates an eGFR using a cystatin-C equation, sometimes combined with creatinine.

\section*{Preparation, Risks, and Limitations}
No fasting is usually required. Venipuncture is the only usual risk. Steroid treatment, thyroid dysfunction, inflammation, smoking, and some other conditions may affect cystatin C independently of filtration. The result should be interpreted with urine albumin, creatinine, trend over time, and the clinical context.

\section*{Understanding Results}
Higher cystatin C generally indicates lower filtration. A persistently reduced eGFR supports chronic kidney disease when present for at least three months or accompanied by other evidence of kidney damage. A single abnormal result does not establish chronic disease.

\section*{References}
National Kidney Foundation guidance on cystatin C and estimated GFR.

\clearpage
\section*{Regulatory Status and Authorization Date}
Regulatory status applies to the specific assay, instrument, manufacturer, and intended use rather than to the test name alone. No single approval date can be assigned to this test category. For a commercial assay, verify the current product in the relevant regulator's database: the U.S. Food and Drug Administration (FDA) clearance, approval, or authorization; India's Central Drugs Standard Control Organisation (CDSCO) licence or permission; the European Union In Vitro Diagnostic Regulation (EU IVDR) CE certificate issued through the applicable conformity-assessment route; or the United Kingdom Medicines and Healthcare products Regulatory Agency (MHRA) registration and UKCA/CE status. Laboratory-developed tests may not have an individual FDA, CDSCO, EU IVDR, or MHRA approval date.
\clearpage
